% !TEX root = main.tex
% -----------------------------------------------------------------------------
% Chapter 3: Methodology
% -----------------------------------------------------------------------------

\chapter{Methodology}
\label{ch:methodology}

\section{Chapter Overview}
This chapter describes the methodology used to evaluate missing-data imputation for a univariate IoT sensor time series. The study follows a controlled experimental design: a clean chronological temperature sequence is selected, synthetic missingness is introduced under several scenarios, multiple baseline imputers are applied to the same masks, and a hybrid residual LSTM-VAE model is evaluated against the same ground truth. The method is hybrid because it combines strong deterministic imputation priors with a neural residual model. This design is important for this dataset because simple interpolation is already very competitive for short random gaps, while boundary gaps require broader seasonal context.

The chapter covers the data preparation process, missingness protocol, baseline methods, proposed model architecture, loss function, training procedure, inference procedure, and evaluation metrics. It also defines the imputation-only runtime measurement used to compare practical computational cost across methods.

\section{Research Design and Workflow}
The research design is experimental and comparative. Every method receives the same observed series and the same missingness mask for a given scenario. Imputation quality is then measured only at the artificially hidden points, where the true values are known. This avoids rewarding a method for simply copying observed values and makes the comparison among methods fair.

The workflow is:
\begin{enumerate}
	\item load the timestamped temperature dataset and sort it chronologically;
	\item construct four missingness scenarios: START, MIDDLE, END, and RANDOM;
	\item build baseline imputations for each scenario;
	\item construct a strong prior imputation for the proposed hybrid model;
	\item train a residual LSTM-VAE on sliding windows of the sequence;
	\item perform windowed imputation and overlap-averaging;
	\item compute MAE, RMSE, visual comparisons, win rate against baselines, and imputation-only runtime.
\end{enumerate}

\begin{figure}[ht]
	\centering
	\placeholderbox{\textbf{Insert Figure: End-to-end methodology pipeline.}\\
	Suggested content: raw temperature series \textrightarrow{} chronological preprocessing \textrightarrow{} missingness masks \textrightarrow{} baseline imputers and hybrid residual LSTM-VAE \textrightarrow{} imputed series \textrightarrow{} MAE/RMSE, plots, and runtime comparison.}
	\caption{End-to-end workflow used in this research, from data ingestion to evaluation.}
	\label{fig:method-pipeline}
\end{figure}

\section{Dataset and Data Selection}
The implementation uses the file \texttt{data/dataset1.csv}, which contains timestamped temperature readings. The datetime column is parsed into a proper time index, the records are sorted chronologically, and the temperature column is extracted as a one-dimensional signal. The final experimental sequence contains $N=26{,}387$ readings.

The study focuses on a single continuous temperature channel. This keeps the imputation task interpretable: the methods are compared on their ability to recover one ordered physical signal rather than on multivariate interactions. Although multivariate sensor fusion is a useful extension, the present work deliberately isolates the univariate imputation problem.

\begin{figure}[ht]
	\centering
	\placeholderbox{\textbf{Insert Figure: Raw temperature time series.}\\
	Suggested content: line plot of temperature against datetime before synthetic masking.}
	\caption{Raw temperature time series after chronological sorting.}
	\label{fig:raw-timeseries}
\end{figure}

\section{Preprocessing and Robust Cleaning}
The raw temperature values are preserved as the ground truth for evaluation. However, the model-fitting process uses a robustly cleaned version of the series to avoid training the neural model to reproduce extreme sensor artefacts. In the implementation, low sentinel-like values below $-20^\circ$C and large local deviations are treated as training artefacts. A centred rolling median and rolling median absolute deviation (MAD) are used:
\[
\mathrm{MAD}_t = \mathrm{median}\left(\lvert x_i - \mathrm{median}(x) \rvert\right),
\]
within a local window. A point is flagged for training cleanup when it is more than a fixed multiple of the robust local scale or falls below the sentinel threshold. These flagged values are filled by time interpolation only for constructing model inputs and training targets.

This distinction is important. The cleaned series stabilises neural training, but the reported MAE and RMSE are still computed against the original raw values at the artificially hidden positions. Therefore, the evaluation remains aligned with the baseline notebooks and does not benefit from changing the test targets.

\section{Missingness Protocol}
Missingness is injected synthetically so that ground truth remains available. The missing rate is fixed at $30\%$ of the sequence. Four scenarios are evaluated:
\begin{itemize}
	\item \textbf{START}: the first $30\%$ of the sequence is hidden;
	\item \textbf{MIDDLE}: a contiguous $30\%$ block in the centre of the sequence is hidden;
	\item \textbf{END}: the final $30\%$ of the sequence is hidden;
	\item \textbf{RANDOM}: $30\%$ of indices are selected uniformly at random using a fixed seed.
\end{itemize}

These scenarios test different imputation difficulties. RANDOM missingness is well suited to interpolation because nearby observations usually remain available. START and END missingness are harder because one side of the temporal context is absent. MIDDLE missingness tests performance on a long internal gap where both past and future context exist, but the gap is large.

\begin{figure}[ht]
	\centering
	\placeholderbox{\textbf{Insert Figure: Missingness scenarios.}\\
	Suggested content: four panels showing START, MIDDLE, END, and RANDOM masks over the same temperature sequence.}
	\caption{Synthetic missingness scenarios used for controlled evaluation.}
	\label{fig:mask-visualisation}
\end{figure}

\section{Baseline Imputation Methods}
Several baseline methods are included to provide interpretable reference points. The same mask is used for each method in a given missingness scenario.

\subsection{Constant Statistical Baselines}
Mean, median, and mode imputation replace every missing value with a single statistic computed from the observed values. These methods ignore temporal structure, but they provide useful lower bounds and highlight whether a method is doing more than estimating the central tendency of the sequence.

\subsection{Sampling Baselines}
Random and hot-deck imputation fill missing positions by sampling values from the observed set. These approaches preserve the marginal distribution better than a single constant, but they do not preserve temporal ordering.

\subsection{Time-Aware Baselines}
Forward fill, last observation carried forward (LOCF), last observation carried backward (LOCB), and linear interpolation exploit temporal order. Forward fill and LOCF use the most recent available reading, with back-filling used to handle initial gaps. LOCB uses the next available reading, with forward filling used when needed. Linear interpolation estimates missing values along straight lines between neighbouring observations and then fills boundary values. These methods are strong baselines for slowly varying physical signals.

\begin{figure}[ht]
	\centering
	\placeholderbox{\textbf{Insert Figure: Baseline comparison.}\\
	Suggested content: zoomed segment showing ground truth, observed values, mean/median, forward fill, and linear interpolation.}
	\caption{Example segment comparing baseline imputations to the ground truth.}
	\label{fig:baseline-visual}
\end{figure}

\section{Hybrid Prior Construction}
The proposed method does not ask the neural network to generate the entire temperature sequence from a zero placeholder. Instead, it first constructs a strong prior estimate and then lets the LSTM-VAE learn a residual correction. This improves stability and makes the model competitive with strong simple baselines.

Two priors are used:
\begin{itemize}
	\item \textbf{Time interpolation prior}: used for MIDDLE and RANDOM missingness, where interpolation has access to useful neighbouring observations.
	\item \textbf{Seasonal KNN prior}: used for START and END missingness, where interpolation lacks one side of the context. The KNN prior matches readings using cyclic calendar features such as time of day and day of year.
\end{itemize}

For each timestamp, calendar features are generated using sine and cosine encodings:
\[
\sin\left(2\pi h_t/24\right),\quad \cos\left(2\pi h_t/24\right),
\]
for hour-of-day, and similarly for day-of-year. A normalised trend feature is also included. These features allow the KNN prior and neural model to use seasonal context without treating calendar time as a discontinuous integer.

\section{Model Input Representation}
All temperature values are standardised using the mean and standard deviation of the observed training values. At each timestep $t$, the model receives a feature vector
\[
\mathbf{u}_t =
[\tilde{x}^{(z)}_t,\; m_t,\; p^{(z)}_t,\; \mathbf{c}_t],
\]
where $\tilde{x}^{(z)}_t$ is the standardised filled value, $m_t$ is the binary observation mask, $p^{(z)}_t$ is the standardised prior value, and $\mathbf{c}_t$ contains the calendar features. Missing entries are filled with the prior value rather than with zero. This prevents the model from seeing artificial discontinuities caused by placeholder values and lets it focus on correcting the prior.

\section{Proposed Method: Residual LSTM-VAE}
The proposed model is a residual LSTM-based Variational Autoencoder. The encoder maps a window of input features into a latent distribution, and the decoder predicts a residual sequence that is added to the prior. Therefore, the final prediction is
\[
\hat{x}^{(z)}_t = p^{(z)}_t + r_t,
\]
where $r_t$ is the residual predicted by the decoder in standardised units.

\subsection{Encoder}
The encoder is an LSTM that processes an input window $\mathbf{u}_{1:T}$. The final hidden state summarises the local temporal context. Two linear layers map this hidden state to a latent mean vector $\boldsymbol{\mu}$ and log-variance vector $\log \boldsymbol{\sigma}^2$:
\[
q_\phi(\mathbf{z} \mid \mathbf{u}_{1:T}) =
\mathcal{N}(\boldsymbol{\mu}, \mathrm{diag}(\boldsymbol{\sigma}^2)).
\]

\subsection{Latent Sampling}
The latent variable is sampled using the reparameterisation trick:
\[
\mathbf{z} = \boldsymbol{\mu} + \boldsymbol{\sigma} \odot \boldsymbol{\epsilon},
\quad
\boldsymbol{\epsilon} \sim \mathcal{N}(\mathbf{0}, \mathbf{I}).
\]
During evaluation, the mean is used to make the reconstruction deterministic.

\subsection{Decoder}
The decoder receives the latent code repeated across the window and concatenated with context features, including the prior and calendar features. A decoder LSTM produces hidden states for each timestep, and a residual head maps those hidden states to scalar residual corrections. The architecture therefore retains the generative latent structure of a VAE while grounding the output in a strong deterministic prior.

\begin{figure}[ht]
	\centering
	\placeholderbox{\textbf{Insert Figure: Residual LSTM-VAE architecture.}\\
	Suggested content: filled value + mask + prior + calendar features \textrightarrow{} encoder LSTM \textrightarrow{} $\mu,\log\sigma^2$ \textrightarrow{} latent code \textrightarrow{} decoder LSTM \textrightarrow{} residual \textrightarrow{} prior + residual.}
	\caption{Architecture of the residual LSTM-VAE used for hybrid time-series imputation.}
	\label{fig:lstm-vae-arch}
\end{figure}

\section{Objective Function}
The model is trained using a VAE-style objective. The reconstruction term is computed only on observed points in the training mask:
\[
\mathcal{L}_{\mathrm{rec}} =
\frac{\sum_{t=1}^{T} m_t \rho(\hat{x}^{(z)}_t - x^{(z)}_t)}
{\sum_{t=1}^{T} m_t + \varepsilon},
\]
where $\rho(\cdot)$ is the Smooth L1 loss. Smooth L1 is used instead of plain mean squared error because it is less sensitive to occasional extreme sensor values.

The KL term regularises the approximate posterior toward the standard normal prior:
\[
\mathcal{L}_{\mathrm{KL}} =
D_{\mathrm{KL}}\left(q_\phi(\mathbf{z}\mid\mathbf{u}) \;\|\; \mathcal{N}(\mathbf{0},\mathbf{I})\right).
\]
A small residual penalty discourages unnecessary deviations from the prior:
\[
\mathcal{L}_{\mathrm{res}} =
\frac{1}{T}\sum_{t=1}^{T}(\hat{x}^{(z)}_t - p^{(z)}_t)^2.
\]
The final objective is
\[
\mathcal{L} =
\mathcal{L}_{\mathrm{rec}}
+ \beta \mathcal{L}_{\mathrm{KL}}
+ \lambda \mathcal{L}_{\mathrm{res}}.
\]
The KL coefficient $\beta$ is warmed up gradually during training. This prevents the KL term from dominating too early and allows the reconstruction pathway to learn a useful representation first.

\section{Training Strategy}
The model is trained with mini-batch gradient descent on sliding windows. The implementation uses windows of length $T=96$, stride $16$, and batch size $128$. The residual LSTM-VAE uses one LSTM layer, hidden dimension $48$, and latent dimension $12$. AdamW is used with learning rate $2\times 10^{-3}$ and weight decay. A ReduceLROnPlateau scheduler lowers the learning rate if the loss stops improving, and gradient clipping is applied to stabilise recurrent training.

Training is performed on masked windows generated from the full sequence. The model sees the filled value, mask, prior, and calendar context, and learns to reconstruct the cleaned standardised target at observed points. The training objective is not computed at artificially missing points, so the model is not directly trained on hidden values that are reserved for evaluation.

\begin{figure}[ht]
	\centering
	\placeholderbox{\textbf{Insert Figure: Training loss curve.}\\
	Suggested content: epoch-wise total loss, reconstruction loss, and learning rate.}
	\caption{Training dynamics of the residual LSTM-VAE under mini-batch optimisation.}
	\label{fig:training-curve}
\end{figure}

\section{Imputation and Inference Procedure}
For each missingness scenario, the appropriate prior is first constructed. The model then performs reconstruction using sliding windows. Each window returns a full residual reconstruction, and overlapping predictions are averaged to produce one reconstructed value per timestamp. This overlap-averaging reduces local variance and avoids relying on a single window boundary.

A conservative residual gate is applied after reconstruction. A small validation mask is created from observed points, and candidate residual blend weights are tested. The neural residual is used only when validation improves both MAE and RMSE over the prior. For MIDDLE and RANDOM missingness, the implementation keeps the prior unchanged because interpolation is already the strongest and most stable choice for those gap structures. For START and END missingness, the residual correction is allowed when the validation check supports it.

The final imputed series preserves all observed raw values and replaces only missing positions with the hybrid model prediction:
\[
x^{\mathrm{imp}}_t =
\begin{cases}
x_t, & m_t=1,\\
\hat{x}_t, & m_t=0.
\end{cases}
\]

\begin{figure}[ht]
	\centering
	\placeholderbox{\textbf{Insert Figure: Windowed reconstruction stitching.}\\
	Suggested content: overlapping windows over the timeline and averaging of predictions into one final imputed series.}
	\caption{Windowed inference and overlap-averaging used to stitch reconstructions into a full-length imputation.}
	\label{fig:window-stitching}
\end{figure}

\section{Evaluation Protocol}
Evaluation is performed only at masked positions. For a given mask, let $\mathcal{M}=\{t:m_t=0\}$ be the set of hidden indices. Mean Absolute Error is computed as
\[
\mathrm{MAE} =
\frac{1}{|\mathcal{M}|}\sum_{t\in\mathcal{M}}
\left|\hat{x}_t - x_t\right|,
\]
and Root Mean Squared Error is computed as
\[
\mathrm{RMSE} =
\sqrt{
\frac{1}{|\mathcal{M}|}\sum_{t\in\mathcal{M}}
\left(\hat{x}_t - x_t\right)^2
}.
\]

In addition to reporting MAE and RMSE for each method, the implementation computes a baseline win rate. For every baseline and every missingness position, the hybrid model receives one win for MAE if its MAE is lower, and one win for RMSE if its RMSE is lower. This gives a compact measure of how often the proposed method improves over the collection of baselines.

\section{Visual and Runtime Comparison}
The evaluation includes qualitative plots for each missingness scenario. For START, MIDDLE, END, and RANDOM, the plotted segment shows the ground truth, observed values, selected baselines, and the hybrid LSTM-VAE imputation. These plots are used to check whether lower error values correspond to visually plausible reconstructions.

Runtime is measured separately from training. The recorded time covers only the imputation procedure: building the imputed output for the method under a fixed mask. For the hybrid model, this includes prior construction, windowed neural reconstruction, overlap averaging, and residual gating. For baselines, it includes only the corresponding baseline imputation operation. This makes the timing comparison focused on deployment-time imputation cost rather than model fitting.

\begin{figure}[ht]
	\centering
	\placeholderbox{\textbf{Insert Figure: Method comparison by missingness scenario.}\\
	Suggested content: four plots comparing ground truth, selected baselines, and hybrid LSTM-VAE for START, MIDDLE, END, and RANDOM.}
	\caption{Qualitative comparison of hybrid LSTM-VAE imputation against baseline methods.}
	\label{fig:qualitative-imputation}
\end{figure}

\begin{figure}[ht]
	\centering
	\placeholderbox{\textbf{Insert Figure: Imputation-only runtime comparison.}\\
	Suggested content: bar chart showing runtime in seconds for the hybrid model and baselines under each missingness scenario.}
	\caption{Imputation-only runtime comparison across methods and missingness scenarios.}
	\label{fig:runtime-comparison}
\end{figure}

\section{Implementation and Reproducibility Details}
All experiments are implemented in Python using NumPy and Pandas for data manipulation, scikit-learn for KNN and baseline utilities, and PyTorch for the residual LSTM-VAE. A fixed random seed is used for missingness generation and sampling baselines. The implementation records the missingness rate, mask type, standardisation statistics, sequence length, stride, batch size, hidden dimension, latent dimension, optimiser settings, and residual blend weights.

After training, the model checkpoint is saved with the model parameters, optimiser state, standardisation parameters, and key architecture settings. This allows the same model configuration to be reused for later evaluation or thesis figure generation.

\section{Chapter Summary}
This chapter presented the methodology used to compare baseline imputation methods with the proposed hybrid residual LSTM-VAE. The method combines deterministic priors with a neural residual model, uses standardised and calendar-aware inputs, evaluates four missingness scenarios, and reports masked-point MAE, RMSE, baseline win rate, visual comparisons, and imputation-only runtime. The next chapter reports and interprets the empirical results produced by this protocol.
